\documentclass[12pt,a4paper]{article}

\usepackage[utf8]{inputenc}
\usepackage[T1]{fontenc}
\usepackage[portuguese]{babel}
\usepackage{amsmath,amssymb,amsfonts}
\usepackage{graphicx}
\usepackage{geometry}
\usepackage{booktabs}
\usepackage{listings}
\usepackage{xcolor}
\usepackage{hyperref}

\geometry{
    top=2.5cm,
    bottom=2.5cm,
    left=2.5cm,
    right=2.5cm
}

\definecolor{codegreen}{rgb}{0,0.6,0}
\definecolor{codegray}{rgb}{0.5,0.5,0.5}
\definecolor{codepurple}{rgb}{0.58,0,0.82}
\definecolor{backcolour}{rgb}{0.96,0.96,0.96}

\lstdefinestyle{pythonstyle}{
    backgroundcolor=\color{backcolour},
    commentstyle=\color{codegreen},
    keywordstyle=\color{magenta},
    numberstyle=\tiny\color{codegray},
    stringstyle=\color{codepurple},
    basicstyle=\ttfamily\footnotesize,
    breakatwhitespace=false,
    breaklines=true,
    captionpos=b,
    keepspaces=true,
    numbers=left,
    numbersep=5pt,
    showspaces=false,
    showstringspaces=false,
    showtabs=false,
    tabsize=4
}

\lstset{style=pythonstyle}

\title{
    \textbf{Trabalho Computacional 04: Simulação de Problemas Eletrostáticos}
    \\[0.3cm]
    \large Modelagem Multifísica de um Eletropropulsor Axissimétrico e Dinâmica de Gotículas
}

\date{}

\begin{document}

\maketitle

\section{Introdução}

O presente trabalho aborda a modelagem e simulação computacional de um propulsor
eletrostático espacial baseado no regime de eletrospray coloidal. O problema é
estruturado em duas etapas complementares:

\begin{enumerate}
    \item \textbf{Parte 1 (Problema Eletrostático Estático):}
    Determinação do campo elétrico crítico de Taylor, cálculo analítico da tensão
    de transição, dedução da formulação fraca axissimétrica para o Método dos
    Elementos Finitos (FEM) e análise teórica do papel do sistema de neutralização
    de carga.

    \item \textbf{Parte 2 (Problema Dinâmico da Gotícula):}
    Modelagem do desacoplamento eletro-hidrodinâmico, determinação das propriedades
    da gotícula (massa e carga), formulação da dinâmica em espaço de estados e
    integração temporal via método de Runge-Kutta para obtenção da velocidade de
    exaustão e do perfil de força ao longo do eixo de simetria.
\end{enumerate}

\section{Parte 1: O Problema Eletrostático Estático}

\subsection{Parâmetros Físicos e Valores Analíticos de Referência}

Considerando os dados fornecidos pelo problema:

\begin{itemize}
    \item Raio do tubo capilar:
    \[
    R = \frac{0{,}12\,\text{mm}}{2}
    = 0{,}06\,\text{mm}
    = 6 \times 10^{-5}\,\text{m}.
    \]

    \item Raio de curvatura inicial do menisco:
    \[
    R_c = R = 6 \times 10^{-5}\,\text{m}.
    \]

    \item Distância intereletrodo:
    \[
    d = 6\,\text{mm} = 6 \times 10^{-3}\,\text{m}.
    \]

    \item Dimensões do anel extrator:
    \[
    R_a = 2\,\text{mm}
    \]
    (raio interno) e
    \[
    R_b = 9\,\text{mm}
    \]
    (raio externo).

    \item Propriedades da formamida:
    tensão superficial
    \[
    \gamma = 0{,}05\,\text{N/m},
    \]
    densidade
    \[
    \rho = 1130\,\text{kg/m}^3.
    \]

    \item Permissividade elétrica do meio (vácuo/ar):
    \[
    \varepsilon_0 \approx 8{,}854187 \times 10^{-12}\,\text{F/m}.
    \]
\end{itemize}

O campo elétrico crítico para o surgimento da instabilidade na superfície líquida
(formação do Cone de Taylor) é dado por:

\begin{equation}
E_{\mathrm{crit}}
=
\sqrt{
\frac{\pi\gamma}
{2\varepsilon_0 R}
}.
\end{equation}

Substituindo os valores:

\begin{equation}
E_{\mathrm{crit}}
=
\sqrt{
\frac{\pi \times 0{,}05}
{2 \times (8{,}854187 \times 10^{-12})
\times (6 \times 10^{-5})}
}
\approx
3{,}845 \times 10^8\,\text{V/m}
\quad
(384{,}5\,\text{MV/m}).
\end{equation}

Aproximando a configuração eletrostática pelo modelo de eletrodo plano infinito
condutor à distância $d$, a diferença de potencial elétrico necessária para
induzir tal campo na ponta é estimada por:

\begin{equation}
V_{\mathrm{analit}}
=
\sqrt{
\frac{\gamma R_c}{\varepsilon_0}
}
\ln\left(
\frac{4d}{R_c}
\right).
\end{equation}

Calculando os termos intervenientes:

\begin{align}
\sqrt{
\frac{\gamma R_c}{\varepsilon_0}
}
&=
\sqrt{
\frac{0{,}05 \times 6 \times 10^{-5}}
{8{,}854187 \times 10^{-12}}
}
\approx
582{,}09\,\text{V},
\\[0.3cm]
\ln\left(
\frac{4d}{R_c}
\right)
&=
\ln\left(
\frac{4 \times 6 \times 10^{-3}}
{6 \times 10^{-5}}
\right)
=
\ln(400)
\approx
5{,}9915.
\end{align}

Resultando na tensão de limiar:

\begin{equation}
V_{\mathrm{analit}}
\approx
582{,}09 \times 5{,}9915
\approx
\mathbf{3487{,}6\,\text{V}}
\quad
(\approx 3{,}49\,\text{kV}).
\end{equation}

\subsection{Dedução da Forma Fraca Axissimétrica (FEM)}

Em coordenadas cilíndricas $(x,r,\theta)$ com simetria azimutal
($\partial/\partial\theta = 0$), a equação de Laplace no domínio meridional
$\Omega$ é expressa por:

\begin{equation}
\nabla \cdot (\varepsilon \nabla \phi)
=
\frac{\partial}{\partial x}
\left(
\varepsilon \frac{\partial\phi}{\partial x}
\right)
+
\frac{1}{r}
\frac{\partial}{\partial r}
\left(
r\varepsilon
\frac{\partial\phi}{\partial r}
\right)
=
0.
\end{equation}

Multiplicando por uma função de teste
$w \in \mathcal{V}$, com

\begin{equation}
\mathcal{V}
=
\left\{
w \in H^1(\Omega)
\mid
w|_{\Gamma_D}=0
\right\},
\end{equation}

e integrando no volume de revolução
\[
dV = 2\pi r\,dr\,dx,
\]
obtém-se:

\begin{equation}
2\pi
\iint_{\Omega}
w
\left[
\frac{\partial}{\partial x}
\left(
\varepsilon
\frac{\partial\phi}{\partial x}
\right)
+
\frac{1}{r}
\frac{\partial}{\partial r}
\left(
r\varepsilon
\frac{\partial\phi}{\partial r}
\right)
\right]
r\,dr\,dx
=
0.
\end{equation}

Integrando por partes no domínio bidimensional $\Omega$:

\begin{equation}
2\pi
\oint_{\partial\Omega}
w\varepsilon
\left(
\nabla\phi\cdot\mathbf{n}
\right)
r\,ds
-
2\pi
\iint_{\Omega}
\varepsilon
\left(
\frac{\partial w}{\partial x}
\frac{\partial\phi}{\partial x}
+
\frac{\partial w}{\partial r}
\frac{\partial\phi}{\partial r}
\right)
r\,dr\,dx
=
0.
\end{equation}

As condições de contorno do sistema são:

\begin{enumerate}
    \item \textbf{Fronteiras de Dirichlet Essenciais:}
    \begin{itemize}
        \item No capilar condutor
        ($x=0$, $0\leq r\leq R$):
        \[
        \phi=V_0
        \implies
        w=0.
        \]

        \item No eletrodo anular aterrado
        ($x=d$, $R_a\leq r\leq R_b$):
        \[
        \phi=0\,\text{V}
        \implies
        w=0.
        \]
    \end{itemize}

    \item \textbf{Eixo de Simetria ($r=0$):}
    A condição natural impõe
    \[
    \left.
    \frac{\partial\phi}{\partial r}
    \right|_{r=0}
    =
    0.
    \]
    O fator $r$ ponderador anula naturalmente o fluxo.

    \item \textbf{Fronteiras Distantes:}
    Condição de truncamento
    \[
    \phi=0\,\text{V}
    \]
    ou Neumann homogêneo
    \[
    \nabla\phi\cdot\mathbf{n}=0.
    \]
\end{enumerate}

Eliminando a constante $2\pi$ e aplicando $w=0$ em $\Gamma_D$, obtém-se a forma
variacional padrão para implementação:

\begin{equation}
\iint_{\Omega}
\varepsilon
\left(
\frac{\partial w}{\partial x}
\frac{\partial\phi}{\partial x}
+
\frac{\partial w}{\partial r}
\frac{\partial\phi}{\partial r}
\right)
r\,dr\,dx
=
0,
\qquad
\forall w\in\mathcal{V}.
\end{equation}

\subsection{Comparação entre a Simulação e a Expressão Analítica}

A formulação analítica clássica adota uma aproximação idealizada de eletrodo
plano maciço infinito. Em contrapartida, o propulsor físico conta com uma abertura
circular de raio $R_a=2\,\text{mm}$ no centro do eletrodo para a passagem do spray.

Essa abertura promove uma alteração do campo elétrico próximo ao eixo axial em
comparação ao caso de uma placa fechada. Em consequência, a tensão numérica
obtida via FEM para alcançar o campo de instabilidade de
$384{,}5\,\text{MV/m}$ pode diferir do valor de $3487{,}6\,\text{V}$ previsto
pelo modelo analítico idealizado.

\subsection{Mecanismo do Sistema Neutralizador}

Devido à polarização do propulsor
($\phi(0)=V_0$ e eletrodo a $0\,\text{V}$), as linhas de campo elétrico apontam
em direção ao anel de extração.

Partículas carregadas positivamente sofrem forças laterais de atração, o que
pode levar à colisão das gotas contra a placa condutora, reduzindo a eficiência
da transferência de quantidade de movimento para o satélite.

Além disso, a exaustão contínua de cargas positivas provocaria o carregamento
elétrico negativo da carcaça do veículo espacial, gerando um potencial adverso
que poderia frear e recolher o feixe de gotículas.

Para contornar tal problema, adota-se um \textbf{cátodo neutralizador}, como
filamentos termoiônicos ou emissores de cátodo oco, posicionado após o orifício
de exaustão. Ele injeta elétrons em quantidade aproximadamente equivalente à
corrente do jato coloidal, reduzindo o carregamento líquido da pluma e permitindo
a operação contínua do sistema de propulsão.



\end{document}